\section{Material and methods}

\begin{frame}{Research objective}
Using Extreme Value Theory (EVT) and copula-based dependence models, this study aims to improve the robustness of tail inference beyond the traditional GEV framework, especially under complex data structures where standard assumptions fail. 
\vspace*{0.3cm}
In order to assess risk measures linking:
\begin{itemize}
\item[-] \textbf{Solar Spectral Irradiance (SSI)};
\item[-] \textbf{Absolute uncertainty of SSI measurements};
\item[-] \textbf{Air density}.
\end{itemize}
\end{frame}


\begin{frame}{Material}

\textbf{Datasets overview}
\vspace*{0.5cm}
\begin{itemize}
  \item \textbf{Solar Spectral Irradiance (SSI)} - NASA SORCE/XPS mission (photodiode 0.1–7 nm, 2005–2019), daily interpolated series [W/m$^2$];
    \vspace*{0.15cm}
  \item \textbf{Absolute Uncertainty (AU) of SSI measurements} - Instrumental 1$\sigma$ uncertainty [W/m$^2$] from the same mission and period for SSI;
    \vspace*{0.15cm}
  \item \textbf{Air Density (AD)} - Obtained from the Brazilian National Meteorological Institute (INMET) station in Brasília, using temperature, pressure, and humidity (Picard et al., 2008).
\end{itemize}
\end{frame}

\begin{frame}{Methods}

\textbf{Modeling steps}
\vspace*{0.5cm}
\begin{enumerate}
  \item \textbf{Univariate modeling}: GEV and BGEV distributions fitted to minima (364 blocks of 15 days) and maxima (90 blocks of 61 days);
  \vspace*{0.15cm}
  \item \textbf{Dependence structure}: parametric bivariate copulas fitted to extreme blocks of SSI, AI, and AD;
  \vspace*{0.15cm}
  \item \textbf{Risk measures}: ``bivariate'' Value at Risk (VaR) and return levels derived from the selected copulas.
\end{enumerate}


\end{frame}